\section{Umbral few-shot, Inverse Scaling y U-shaped Scaling}

Con la definición operativa en mano, el siguiente paso no es buscar de inmediato un mecanismo misterioso, sino aprender a leer las figuras. El few-shot prompting coloca varios exemplar de input--output en el contexto para que un frozen language model siga prediciendo la respuesta; si el modelo pequeño solo adivina y el grande puede aprovechar los exemplar para resolver la tarea, se tiene lo que en clase se llama few-shot emergence.

\subsection{Cómo el few-shot prompting se convierte en un experimento medible}

El experimento requiere definir con claridad el prompt, el input no visto, la target label y el chance baseline. Si la exactitud aleatoria de una tarea de clasificación binaria es $0.5$, que un modelo pequeño quede alrededor de $0.5$ no prueba que «no entienda nada», pero al menos indica que la interface actual no convierte su conocimiento en un desempeño utilizable; solo cuando el modelo grande supera de forma sostenida el baseline puede decirse que la capacidad es medible bajo esa interface. Dicho de otro modo, juzgar una capacidad no es la impresión subjetiva que deja leer una salida del modelo, sino aplicar el mismo elicitation procedure a una serie de scale checkpoints, compararlos y comprobar si la mejora rebasa lo que pueden explicar el ruido de muestreo, el desbalance de clases y la elección fortuita del prompt.

\lecturefigure{slide-07.jpg}{Definición de emergencia en una few-shot prompted task}{Diapositivas oficiales, p. 7; Wei et al., 2022}

Para una clasificación balanceada de $K$ clases, puede usarse como apoyo para comparar el chance-normalized gain:
\[
G_{\mathrm{chance}}=\frac{a-1/K}{1-1/K},
\]
donde $a$ es la exactitud y $1/K$ es la línea base aleatoria. $G_{\mathrm{chance}}=0$ significa que no se supera el azar y $G_{\mathrm{chance}}=1$ significa la calificación perfecta. Facilita observar «cuánto se supera el chance» entre tareas con distinto número de clases, pero no elimina la prompt sensitivity ni la varianza de las muestras.

\lecturefigure{slide-08.jpg}{Curvas de few-shot emergence en varias tareas}{Diapositivas oficiales, p. 8; Wei et al., 2022}

\begin{knowledgebox}{Lectura de la figura: primero el baseline, luego el punto de inflexión}
En cada gráfica pequeña, el eje horizontal es el número de FLOPs de entrenamiento o el tamaño del modelo, y el vertical es la metric de cierta tarea. Primero, ubica el random baseline; segundo, observa si los puntos de los modelos pequeños oscilan alrededor del baseline; tercero, determina si el modelo grande supera el baseline en varios puntos consecutivos, y no solo en un punto aislado que podría ser ruido; cuarto, compara si el punto de inflexión coincide entre distintas familias de modelos. Lo «repentino» describe una curva de tarea con muestreo escaso; no significa que la representación interna se genere instantáneamente entre dos checkpoint contiguos.
\end{knowledgebox}

\teachervoice{El ponente explica estas figuras eje por eje: el horizontal es la escala, el vertical es el desempeño, y la línea punteada o el valor conocido es el nivel aleatorio. Lo que se enfatiza en clase no es «ver la curvatura a simple vista», sino el criterio de que el modelo pequeño no supera el azar y el grande empieza a superarlo de forma estable.}

\subsection{Dos tareas concretas: pruebas de conocimiento y transcripción fonética}

Las gráficas pequeñas y abstractas hacen que el lector pase por alto las task mechanics, por lo que las diapositivas presentan a continuación dos ejemplos más concretos. Uno es un conjunto de preguntas de opción múltiple de varias disciplinas; el otro consiste en convertir oraciones en inglés al International Phonetic Alphabet (IPA, Alfabeto Fonético Internacional). Ambos exigen que el modelo transforme los ejemplos del contexto en un nuevo formato de salida.

\lecturefigure{slide-09.jpg}{Fenómeno de umbral few-shot en preguntas de opción múltiple de varias tareas}{Diapositivas oficiales, p. 9; datos de Hendrycks et al., 2020}

En este tipo de benchmark, la raw accuracy mezcla conocimiento, comprensión del enunciado, comparación de opciones y formato de salida. Si el modelo pequeño no ejecuta de forma estable cualquiera de estos pasos, la calificación total queda cerca del azar; el aumento de escala puede hacer que varios pasos se vuelvan «suficientemente confiables» al mismo tiempo, lo que produce un punto de inflexión notorio en la aggregate metric. La figura por sí sola no nos dice qué paso mejora primero.

\lecturefigure{slide-10.jpg}{Comportamiento emergente en la tarea de transcripción del inglés al IPA}{Diapositivas oficiales, p. 10; BIG-Bench}

El ejemplo del IPA muestra un espacio de salida de baja frecuencia: la entrada es una oración común en inglés, pero el objetivo contiene secuencias de símbolos fonéticos que rara vez se le pide al modelo generar conforme a reglas. Al leer la figura conviene separar «si reconoce la intención de la tarea», «si domina la correspondencia de fonemas» y «si puede producir caracteres especiales de forma estable»; que la exactitud de la tarea cruce el umbral solo indica que la cadena completa por fin es utilizable, lo cual no equivale a que cada subcapacidad nazca al mismo tiempo.

\subsection{Por qué la inverse scaling puede ser solo la primera mitad de la curva}

Si solo se observa que el small model supera al medium model, es fácil concluir que «seguir aumentando la escala lo hará cada vez peor». En clase se usa un ejemplo deliberately provocative de selección de palabras para mostrar que, al crecer, el modelo puede reproducir con mayor fidelidad las asociaciones comunes de la distribución de entrenamiento y, por ello, alejarse primero del objetivo de la tarea; al seguir creciendo, una instruction/context understanding más sólida puede revertir de nuevo la tendencia.

\lecturefigure{slide-11.jpg}{Inverse scaling can become U-shaped}{Diapositivas oficiales, p. 11; Wei, Tay, and Le, 2022}

\begin{knowledgebox}{Lectura de la figura: tres tramos para explicar la U-shape}
El small model, el medium model y el large model de la figura no deberían recibir solo la etiqueta de «bueno/malo». En el primer tramo, es posible que el modelo aún no haya aprendido un prior fuerte; en el segundo, el prior se vuelve fuerte, pero el control de la tarea es insuficiente; en el tercero, la context sensitivity, el instruction following o capacidades de composición de nivel superior alcanzan al prior. La verdadera conclusión didáctica es que, mientras la curva no cubra escalas mayores, no debe extrapolarse una pendiente negativa local como si fuera una ley permanente.
\end{knowledgebox}

\begin{warningbox}{Los checkpoint escasos generan dramatismo visual}
Si el eje horizontal tiene solo tres o cuatro modelos, el punto mínimo de la U-shape y el «umbral» solo pueden ubicarse de manera aproximada. Distintos prompt, seed, metric o particiones de datos también pueden desplazar la curva. La formulación correcta es «se observó una inversión de tendencia en la familia de modelos y la configuración evaluadas», y no «toda inverse scaling acabará resolviéndose con la escala».
\end{warningbox}

\subsection{Emergent prompting technique: el método también tiene un umbral}

Lo anterior trata de la task ability; después, las diapositivas desplazan la perspectiva hacia la intervention: un mismo método de RLHF o de prompting puede no aportar ganancia positiva en un modelo pequeño y mostrar delta solo en uno grande. Por lo tanto, la eficacia de un método no puede validarse únicamente en un modelo pequeño y barato para luego extrapolarse sin más.

\lecturefigure{slide-12.jpg}{Una intervención puede tener ganancia negativa en un modelo pequeño y volverse positiva en uno grande}{Diapositivas oficiales, p. 12; Wei et al., 2022}

\begin{knowledgebox}{Lectura de la figura: se compara el delta dentro de una misma escala}
Primero compara el baseline con la intervention en cada escala, en lugar de comparar solo las puntuaciones absolutas de los extremos izquierdo y derecho. En el modelo pequeño, la línea naranja queda por debajo de la azul, lo que indica que la intervención ocupa capacidades de representación u optimización que el modelo aún no tiene; en el modelo grande, la línea naranja supera a la azul, lo que indica que la misma intervención por fin puede aprovechar una base capability más fuerte. Este ejemplo también recuerda que hay que distinguir entre la escala de preentrenamiento, el prompting y el weight-changing post-training.
\end{knowledgebox}

\teachervoice{Un estudiante pregunta si todos estos modelos pasaron por instruction fine-tuning, y el ponente responde con claridad que este conjunto de figuras de base-model emergence no aplicó instruction fine-tuning. Después usa el RLHF para explicar que probar un método solo en modelos pequeños y observar que falla no basta para demostrar que tampoco funcione en modelos más grandes.}

\subsection{Resumen de la sección}

La few-shot emergence es un fenómeno a nivel de tarea relativo al chance baseline; la tendencia negativa local de la inverse scaling puede revertirse a mayor escala; y una intervention de prompting o de post-training también puede tener su propio umbral de capacidad. Los tres puntos se oponen a un error común: extrapolar linealmente a todo el intervalo de escalas a partir de experimentos con un número limitado de modelos pequeños.
